% Generated draft: requires mathematical and source review.
\documentclass[11pt]{article}
\usepackage[margin=25mm]{geometry}
\usepackage{fontspec}
\setmainfont{Latin Modern Roman}
\usepackage{amsmath,amssymb,mathtools,unicode-math}
\setmathfont{Latin Modern Math}
\usepackage{enumitem,needspace,xurl,hyperref}
\setlength{\parindent}{0pt}
\setlength{\parskip}{5pt}
\setlength{\emergencystretch}{3em}
\newcommand{\N}{\mathbb{N}}
\newcommand{\Z}{\mathbb{Z}}
\newcommand{\Q}{\mathbb{Q}}
\newcommand{\R}{\mathbb{R}}
\newcommand{\C}{\mathbb{C}}
\newcommand{\F}{\mathbb{F}}
\newcommand{\A}{\mathbb{A}}
\newcommand{\PP}{\mathbb{P}}
\newcommand{\E}{\mathbb{E}}
\newcommand{\eps}{\varepsilon}
\newcommand{\dd}{\,\mathrm{d}}
\begin{document}
\section*{UnsolvedMath ID: 30004964}
\textbf{Source problem code:} catalogue entry
\subsection{Definitions and mathematical statement}
Let \(\mathcal G\) be a graph class. A graph class is \emph{monotone} if it is closed under taking subgraphs, including subgraphs obtained by deleting vertices or edges.

Under the usual convention of finite simple undirected graphs, an \(r\)-shallow minor of a graph \(G\), for \(r\in\Z_{\geq 0}\), is obtained from a subgraph of \(G\) by contracting pairwise vertex-disjoint connected branch sets, each having radius at most \(r\) within that branch set, and discarding loops and parallel edges. Under this convention, \(\mathcal G\) has \emph{bounded expansion} if
\[
\forall r\in\Z_{\geq 0}\ \exists b_r\in\R_{\geq 0}\ \forall G\in\mathcal G\ \forall H,
\qquad
\bigl(H\text{ is an }r\text{-shallow minor of }G,\ |V(H)|>0\bigr)
\Longrightarrow
\frac{2|E(H)|}{|V(H)|}\leq b_r.
\]
The bound \(b_r\) may depend on \(r\) and the class \(\mathcal G\), but not on \(G\) or \(H\). The supplied record does not explicitly specify the graph domain; the convention used in this explanatory definition requires confirmation against the report.

\textbf{Conjecture 5, as identified by the supplied curation.} If a monotone graph class \(\mathcal G\) is \((C,s)\)-sketchable for some real \(C>1\) and some admissible parameter \(s\), then \(\mathcal G\) has bounded expansion. The parameters are existentially quantified for the whole class.

\textbf{Unresolved qualification.} The supplied record gives neither the definition of \((C,s)\)-sketchability nor the admissible domain of \(s\). Consequently, the preceding conjecture preserves the supplied target but is not a self-contained mathematical formulation: the meaning of its sketchability hypothesis must be recovered from the cited report before mathematical review or use.

The supplied contextual passage says that Conjecture 5 would follow from a conjecture attributed to Thomassen, described qualitatively as asserting that every graph of sufficiently large average degree contains a subgraph of large girth and large average degree. This contextual description supplies neither the precise quantifiers nor the dependence of the required degree threshold on the target parameters; it is retained as source context rather than as a separately formalized assertion.
\subsection{Short English statement}
The supplied conjecture asks whether a graph class closed under subgraphs must have bounded expansion whenever it is \((C,s)\)-sketchable for some \(C>1\) and some \(s\). The record omits the definition of sketchability and the domain of \(s\), so the hypothesis requires clarification from the cited report. The supplied open status and claimed connection to Thomassen's conjecture are dated source assessments, not independently verified conclusions.
\subsection{Sources}
\begin{itemize}
\item Supplied local UnsolvedMath registry record: permanent ID \(30004964\), exact source problem code \texttt{catalogue entry}, title \emph{Sketchability and Bounded Expansion of Graph Classes}, collection \emph{Oberwolfach Reports Open Problems}. No individual proposer is identified; the supplied proposed year is 2022.
\item Supplied source citation: \emph{Graph Theory (2023)}, Oberwolfach Reports, DOI \texttt{10.4171/owr/2022/1}, pp.\ 5--78. Supplied source URL: \url{https://doi.org/10.4171/owr/2022/1}. The background describes these reports as documenting mathematical workshops and their open-problem discussions.
\item Supplied literature source, labelled ``original OWR report'': \url{https://doi.org/10.4171/OWR/2022/1}. The curation identifies the extracted question as Conjecture 5. The source record labels the statement \texttt{corrected\_verified} and reports that it was checked against the official report; this verification claim is supplied metadata and has not been independently checked here.
\item Supplied literature source, labelled ``Thomassen Girth in graphs'': \url{https://doi.org/10.1016/0095-8956(83)90067-9}. The contextual excerpt attributes the related conjecture to Thomassen and cites it as reference [3]; no complete reference-[3] bibliographic entry is supplied.
\item Literature assessment checked on \textbf{2026-08-22}: the supplied curation labels this problem open, says that the report presents Conjecture 5 as unresolved, and describes the related Thomassen conjecture as still open. This is supplied dated curation, not independently verified literature assessment. The background explicitly requests independent verification of time-sensitive status and cited claims.
\end{itemize}
\end{document}
